\personalAcknowledgementEnglish{
First and foremost, I want to thank my advisor, Mark Silberstein. You kept pushing me, and at the same time you trusted my work at every step along the way. Your feedback was always constructive and genuinely helpful, and, importantly, you taught me how much it matters to keep questioning my own arguments, to fail fast, and then to keep pushing anyway. These are lessons I will carry with me, and I know they will make me a better researcher and a better engineer. Once, in a struggle I had, you told me something I carried with me: ``If everything had been easy for you up until now, it only meant you had never faced a barrier hard enough to cross''. In this research we managed to cross a few of those. Thanks to you I gained skills I value deeply, technical and interpersonal alike. I had the chance to do things I never imagined I would, from studying such a fascinating topic in depth, to reverse engineering real hardware, to finishing a thorough thesis, a research I care about, and a paper. Thank you for guiding me through all of it.

To my parents, thank you for pushing me forward even when the frustration ran high. Dad, thank you for advising me all the way, both morally and technically, and Mom, thank you for keeping me on track and helping me hold everything in proportion.

Ella, Matan, Emma and Liam, thank you. You are my joy and I love you. Some of the best moments were simply the time we spent together, playing, our Shabbat dinners, and the holidays. Even on the tightest deadline, I could never resist the temptation to come and play with you. None of this would have been possible without you.

To my friends, thank you. Yaheli, thank you for listening through every difficulty and joy in every achievement. Yael, thank you for always helping me put things back into proportion. Jonatan, Selam, and Ofek, thank you for always listening to my ideas, giving me honest feedback, celebrating with me when an idea worked, and pushing me to look for the next one when it did not. Thank you for enduring hours upon hours of me insisting that the sauna is the best place in the world to have ideas, and for staying with me through the whole journey.

Finally, to all the members of the ACSL group, who are also my friends, thank you. A special thanks to Jakob, Michael, Awatef, Daniel, and Noga. Our conversations, about research, but just as often about philosophy, politics, travels across the world, and even learning Chinese, were some of the most fun I had. A few of those conversations should never leave that room, and that is probably for the best. This would not have been possible without you.
}

\personalAcknowledgementHebrew{
ראשית, אני רוצה להודות למנחה שלי, מרק זילברשטיין. דחפת אותי קדימה והאמנת במחקר שלי לכל אורך הדרך. בשיחות שלנו לימדת אותי מהו פידבק איכותי וכנה. לימדת אותי לפעול בצורה מתודולוגית ועקבית. לימדת אותי כמה חשוב לערער על המסקנות שאני מגיע אליהן, להיכשל מהר ולהמשיך לתקוף את הבעיות מכיוונים אחרים. בזכותך נהייתי מהנדס טוב יותר וחוקר טוב יותר. במהלך אחד הקשיים שנתקלתי בהם במחקר אמרת: ``אם עד עכשיו הכל היה לך יחסית קל, אז זה סימן שלא נתקלת באתגר קשה מספיק''. במחקר הזה עברנו כמה אתגרים שלא האמנתי שאעבור. בזכותך רכשתי מיומנויות חשובות ועמוקות, טכניות ובינאישיות. זכיתי לעשות דברים שלא העליתי בדמיוני שתהיה לי הזדמנות לעשות, החל מללמוד נושא כה מרתק, דרך הנדסה לאחור של חומרה, ועד כתיבת תזה, מחקר שאכפת לי ממנו ואפילו מאמר. תודה שליווית אותי בכל זה.

להוריי, תודה שהמשכתם לדחוף אותי קדימה גם בתקופות קשות ומתסכלות. אבא, תודה שתמיד שמחת לתרום ולעזור בפן הטכני והאישי, ואמא, תודה ששמרת עליי שפוי ותמיד עזרת לי לשמור על פרופורציות.

אלה, מתן ושני האחיינים המתוקים שלי, אמה וליאם, תודה. אני אוהב אתכם. חלק מהרגעים הטובים ביותר בשנים האלה היו פשוט הזמן שבילינו יחד, במשחק, בארוחות השבת ובחגים. גם בדדליין הכי צפוף לא יכולתי לעמוד בפיתוי לבוא ולשחק אתכם. כל זה לא היה אפשרי בלעדיכם.

תודה גם לחברים שלי. יהלי, תודה שהקשבת לכל קושי ושמחת בכל הצלחה. יעל, תודה שעזרת לי לשמור על פרופורציות לאורך הדרך. יונתן, סלאם ואופק, תודה שתמיד הקשבתם לרעיונות שלי, נתתם לי פידבק, שמחתם איתי כשהם עבדו, ודחפתם אותי לחפש את הרעיון הבא כשלא. וכמובן תודה שעמדתם בשעות על גבי שעות שבהן התעקשתי שסאונה היא המקום הטוב בעולם להעלות בו רעיונות, ושהייתם לצידי לאורך כל הדרך.

לבסוף, לכל חברי \textenglish{ACSL}, שכמובן הם לא פחות מחברים שלי, תודה. תודה מיוחדת ליעקב (ג'ייקוב), מיכאל, עואטף, דניאל ונוגה. השיחות שלנו, על מחקר, אבל לא פחות מכך על פילוסופיה, פוליטיקה, טיולים ברחבי העולם, ואפילו ללמוד סינית, היו מהדברים הכי כיפיים שחוויתי. כנראה חלק מהשיחות שלנו לא צריכות לצאת מאותו משרד, וטוב שכך. כל זה לא היה אפשרי בלעדיכם.
}
